\begin{textsample}{POS Dim 4 – 2016\_c5 – Score 2.00 – 2016\_c5/enade\_2016\_med\_vet\_07e2ff493e.txt}  \label{ex:f4_pos_105}
CAPÍTULO 5 – DISTRIBUIÇÃO DOS CONCEITOS 5.1 - PANORAMA NACIONAL DA DISTRIBUIÇÃO DOS CONCEITOS O Conceito Enade é um indicador de qualidade que avalia o desempenho dos estudantes a partir dos resultados do Enade.
Os valores possíveis vão de 1 ( pior situação ) a 5 ( melhor situação ).
As Unidades de \textbf{Observação} com apenas um ou sem nenhum concluinte participante não obtêm o Conceito Enade, ficando Sem Conceito ( SC ).
A Tabela 5.1 apresenta a quantidade e distribuição de cursos de Medicina Veterinária participantes do Enade/2016, por faixa de conceito e Grande Região.
Enfatiza -se, mais uma vez, que a diferença entre os cursos tabulados neste capítulo e no capítulo 2 corresponde aos cursos sem conceito, ou seja àqueles, em princípio, sem alunos concluintes que houvessem participado da prova ou àqueles que tivessem tão somente um aluno concluinte19.
Observando -se os dados da Tabela 5.1, nota -se que, dos 179 cursos participantes, 71 ( 39,7 \% ) classificaram -se com conceito 3, o valor modal.
Esse foi também o conceito modal em três regiões : Norte ( 54,5 \% ), Sudeste ( 49,4 \% ) e Sul ( 38,6 \% ).
Nas demais regiões, a moda foi o conceito 4, atribuído para 56,5 \% dos cursos da região Nordeste e para 45,0 \% dos cursos da região Centro-Oeste.
O conceito 4 foi o segundo mais frequente em nível nacional ( 30,7 \%, correspondendo a 55 cursos ) e o conceito 2, o terceiro ( 18,4 \%, correspondendo a 33 cursos ).
Houve, ainda, 16 cursos ( 8,9 \% ) que receberam conceito 5 e quatro cursos ( 2,2 \% ) que receberam conceito 1.
Dos 179 cursos de Medicina Veterinária oferecidos no Brasil, nenhum curso ficou sem conceito ( SC ). 19 Estes cursos com somente um concluinte não foram considerados no capítulo 2.
A região Norte participou com 11 cursos ou 6,1 \% do total nacional.
Desses, seis cursos receberam o conceito 3, o que equivale a 54,5 \% do total regional, além de ser o conceito modal na região, como já comentado.
A três cursos ( 27,3 \% ) atribuiu -se conceito 2 e a dois cursos ( 18,2 \% ), conceito 4.
Nenhum curso ficou sem conceito ( SC ) ou com conceito 1 ou 5 nessa região.
A região Nordeste participou com 23 cursos ou 12,8 \% do total nacional.
Nessa região, 13 cursos ( 56,5 \% em termos regionais ) obtiveram conceito 4, o conceito modal para a região, como já comentado.
Os conceitos 1 e 2 foram atribuídos, respectivamente, a um e a quatro cursos ( 4,3 \% e 17,4 \%, respectivamente ).
Ainda nessa região, cinco cursos ( 21,7 \% ) foram avaliados com o conceito 3.
Nenhum curso ficou sem conceito ( SC ) ou recebeu conceito 5 nessa região.
Dos 81 cursos participantes da região Sudeste ( 45,3 \% do total nacional ), o conceito 3 foi o conceito modal, como já anteriormente assinalado, tendo sido obtido por 40 cursos ( 49,4 \% ).
O conceito 2 foi atribuído a 16 cursos ( 19,8 \% ) e o conceito 4, a 15 cursos ( 18,5 \% ).
Além disso, dez cursos ( 12,3 \% ) receberam o conceito 5.
Nessa região, nenhum curso ficou sem conceito ( SC ) ou com conceito 1.
Os 44 cursos da região Sul corresponderam a 24,6 \% do total nacional.
Como já indicado, houve predominância do conceito 3, atribuído a 38,6 \% dos cursos da região ( conceito modal ), ou seja, atribuído a 17 dos 44 cursos participantes na região Sul.
Os conceitos 2 e 4 foram atribuídos a, respectivamente, cinco e 16 cursos.
Seis cursos receberam conceito 5.
Nenhum curso ficou sem conceito ( SC ) ou recebeu conceito 1.
Já dos 20 cursos participantes na região Centro-Oeste ( 11,2 \% do total nacional ), nove receberam conceito 4, o conceito modal ( 45,0 \% ), como já fora destacado.
A cinco cursos foi atribuído conceito 2, e a seis, conceitos 1 e 3 ( três cursos, cada ).
Nessa região, nenhum curso ficou sem conceito ( SC ) ou com conceito 5.
O Gráfico 5.1 apresenta a distribuição do Conceito Enade segundo Grande Região : poligonais mais à esquerda representam Grandes Regiões com pior distribuição, e poligonais mais à direita, Grandes Regiões com melhores conceitos.
A região Norte ( linha azul ) apresenta os cursos com conceitos mais baixos, ao passo que as regiões Nordeste ( linha verde ) e Sul ( linha roxa ) apresentam as distribuições com valores maiores.
Cursos em IES situadas nas regiões Sudeste ( linha vermelha ) e Centro-Oeste ( linha preta ) apresentam uma situação intermediária. 5.2 - CONCEITOS POR CATEGORIA ADMINISTRATIVA E GRANDE REGIÃO A Tabela 5.2 apresenta a distribuição dos cursos participantes do Enade/2016 de Medicina Veterinária, por Categoria Administrativa, de acordo com os conceitos por eles alcançados, segundo as Grandes Regiões brasileiras.
Dos 179 cursos participantes, 61 ( 34,1 \% ) eram ministrados em Instituições Públicas e 118 ( 65,9 \% ), em Instituições Privadas.
De acordo com as informações da Tabela 5.2, em termos nacionais, entre Instituições Públicas e Privadas, observa -se que, dos 16 cursos avaliados com conceito 5, 14 eram oferecidos em IES Públicas e duas, em IES Privadas.
Dos 61 cursos participantes de IES Públicas, o conceito 4 foi o valor modal, atribuído a 32 cursos ( 52,5 \% ).
Entre os demais cursos de Instituições Públicas participantes, dois cursos ( 3,3 \% ) foram avaliados com conceito 1, quatro cursos ( 6,6 \% da categoria ) obtiveram conceito 2, nove cursos ( 14,8 \% ) receberam conceito 3 e, como já comentado, 14 cursos ( 23,0 \% ) receberam conceito 5.
Nesta categoria, nenhum curso ficou sem conceito ( SC ).
Na rede privada, o conceito modal foi 3, atribuído a 62 cursos, o correspondente a 52,5 \% dos 118 cursos da categoria.
Entre os demais cursos participantes, dois ( 1,7 \% ) receberam conceito 1, e 29 ( 24,6 \% ), conceito 2.
O conceito 4 foi atribuído a 23 cursos ( 19,5 \% ) e o conceito 5, a dois cursos ( 1,7 \% ).
Nesta Categoria Administrativa, nenhum curso ficou sem conceito ( SC ).
Na análise por região, observa -se que, na região Norte, as Instituições Privadas participaram com seis cursos ( 54,5 \% do total regional ), dos quais nenhum ficou sem conceito ( SC ) ou recebeu conceito 1 ou 5.
O conceito modal para as Instituições Privadas na região foi 2, atribuído a três cursos, correspondendo a 50,0 \% dessa categoria na região.
Dois cursos foram avaliados com conceito 3 e um curso, com conceito 4.
As Instituições Públicas participaram com cinco cursos na região Norte ( 45,5 \% do total regional ), dos quais quatro obtiveram conceito 3, o valor modal.
O outro Na região Nordeste, a rede privada concentrou seis dos 23 cursos participantes, o equivalente a 26,1 \% do total da região, a menor das proporções desta categoria entre as regiões.
Dentre os cursos oferecidos por IES Privadas no Nordeste, três cursos foram avaliados com conceito 2 e outros três com conceito 3.
As Instituições Públicas da região Nordeste participaram com 17 cursos ( 73,9 \% do total da região ), a maior das proporções desta categoria nas regiões.
Desses, 13 obtiveram conceito 4, o valor modal.
Os demais cursos receberam conceito 3 ( dois cursos ) e conceitos 1 e 2 ( um curso, cada ).
Nenhum curso ficou sem conceito ( SC ) ou com conceito 5.
Na região Sudeste, a proporção de cursos da rede privada, 82,7 \%, foi mais elevada do que a registrada nas demais regiões brasileiras, correspondendo a 67 dos 81 cursos participantes.
Nessa categoria e região, o conceito modal foi 3, atribuído a 38 cursos.
Nenhum curso ficou sem conceito ( SC ) ou conceito 1.
Os demais foram avaliados com conceito 2 ( 15 cursos ), conceito 4 ( 12 cursos ) e conceito 5 ( dois cursos ).
Entre os 14 cursos oferecidos por Instituições Públicas na região Sudeste, o conceito modal foi 5, atribuído a oito cursos.
Os demais cursos receberam os conceitos 2 ( um curso ), 3 ( dois cursos ) e 4 ( três cursos ).
Na região Sudeste, nesta categoria, nenhum curso ficou sem conceito ( SC ) ou recebeu conceito 1.
As Instituições Privadas concentraram 27 dos 44 cursos participantes da região Sul, 61,4 \% do total regional.
Desses, 16 obtiveram conceito 3, o conceito modal.
Nesta combinação de Categoria Administrativa e Grande Região, os conceitos 2 e 4 foram atribuídos a quatro e a sete cursos, respectivamente.
Nenhum curso ficou sem conceito ( SC ) ou recebeu conceito 1 ou 5.
As Instituições Públicas na região Sul participaram com 17 cursos, 38,6 \% dos cursos da região.
O conceito modal foi 4, atribuído a nove cursos.
Aos demais cursos foram atribuídos os conceitos 2 e 3 ( um curso, cada ) e 5 ( seis cursos ).
Nenhum curso da região Sul oferecido por IES Pública ficou sem conceito ( SC ) ou recebeu os conceitos 1.
Na região Centro-Oeste, 12 dos 20 cursos participantes eram de Instituições Privadas ( 60,0 \% em termos regionais ).
Desses, quatro receberam conceito 2, o conceito modal.
Entre os demais, dois receberam conceito 1, seis receberam conceitos 3 e 4 ( três cursos ).
Nenhum curso ficou sem conceito ( SC ) ou com conceito 5.
Quanto aos oito cursos oferecidos por Instituições Públicas na região Centro-Oeste ( 40,0 \% do total regional ), nota -se seis cursos ficaram com conceito 4, conceito modal.
Dois cursos

% matched lemmas: observação
\end{textsample}
